\documentclass{article}
\usepackage[utf8]{inputenc}
\usepackage{amsmath,amssymb}
\usepackage[most]{tcolorbox}
\usepackage{geometry}
\geometry{margin=2.5cm}

\newcommand{\step}[1]{\par\noindent\hangindent=3.5em\hangafter=1%
  \makebox[3.5em][l]{\textbf{#1.}}}
\newcommand{\steptag}[1]{\hfill{\small\textsf{[#1]}}}

\newtcolorbox{proofbox}[1]{
  colback=gray!5, colframe=gray!60,
  fonttitle=\bfseries, title={#1},
  breakable, enhanced,
  left=4pt, right=4pt, top=4pt, bottom=4pt,
  before skip=10pt, after skip=10pt
}

\begin{document}

\begin{proofbox}{Parameterized complete quotient package: for every univer...}
\step{1} Parameterized complete quotient package: for every universal def-stage1-polar-witness-data tuple W=(C\_rect,C\_ch,C\_pol,C\_grp,C\_path,C\_der,e\_rect,kappa\_ch,kappa\_pol,kappa\_der,delta\_*,epsilon\_*\textasciicircum{}r,e\_S1,r\_iso) with C\_rect,C\_ch,C\_pol,C\_grp,C\_path,C\_der>=1, 0<e\_rect<=1/C\_rect, 0<kappa\_ch,kappa\_pol,kappa\_der<=1/2, delta\_*=min\{1/4,kappa\_ch/(4*C\_ch),kappa\_pol/(4*C\_pol)\}, epsilon\_*\textasciicircum{}r=min\{1/4,kappa\_ch/(4*C\_ch),kappa\_pol/(4*C\_pol),kappa\_der/(8*C\_der),1/C\_grp,delta\_*/(12*C\_path*C\_grp)\}, e\_S1=min\{e\_rect,epsilon\_*\textasciicircum{}r/C\_rect\}, and r\_iso=min\{delta\_*/4,kappa\_der/(8*C\_der)\}, there exist universal e\_quot\textasciicircum{}r>0 and epsilon\_B\textasciicircum{}r>0 with epsilon\_B\textasciicircum{}r=min\{epsilon\_*\textasciicircum{}r,e\_quot\textasciicircum{}r\} such that, for every finite-dimensional exact-unit epsilon\_r-C*-algebra with 0<=epsilon\_r<=epsilon\_B\textasciicircum{}r and 1<dim\_C calX<infinity, and for displayed data satisfying all of the following: (A\_2) for every V in calUbar\_delta\_* and A\textasciicircum{}par in B\_\{2delta\_*\}\textasciicircum{}\{icalH\}(0) there is a unique g\_V(A\textasciicircum{}par) in B\_\{2delta\_*\}\textasciicircum{}\{calH\}(0) with f\_V(A\textasciicircum{}par+g\_V(A\textasciicircum{}par))=0, where f\_V(A)=(1/2)*(((J+A\textasciicircum{}dagger) bold-dot V\textasciicircum{}dagger) bold-dot (V bold-dot (J+A))-J), chi\_V(A\textasciicircum{}par)=V bold-dot (J+A\textasciicircum{}par+g\_V(A\textasciicircum{}par)) lies in calU, ||Dg\_V(A\textasciicircum{}par)||<=C\_ch*(epsilon\_r+delta\_*), ||D\_\{A\textasciicircum{}perp\}f\_V(A\textasciicircum{}par+g\_V(A\textasciicircum{}par))-I\_calH||<=C\_ch*(epsilon\_r+delta\_*)<1, and these C\textasciicircum{}1 charts cover calU; (A\_4) set Pi\_delta\_*:calU x B\_\{delta\_*\}\textasciicircum{}\{calH\}(J)->calX by Pi\_delta\_*(U,K):=U bold-dot K and set S\_delta\_*:=Pi\_delta\_*(calU x B\_\{delta\_*\}\textasciicircum{}\{calH\}(J)); Pi\_delta\_* is a C\textasciicircum{}1 diffeomorphism onto S\_delta\_* with unique inverse maps u\_delta:S\_delta\_*->calU and h\_delta:S\_delta\_*->B\_\{delta\_*\}\textasciicircum{}\{calH\}(J) satisfying X=u\_delta(X) bold-dot h\_delta(X) for every X in S\_delta\_* and u\_delta(U)=U and h\_delta(U)=J for every U in calU; (A\_5) the displayed maps mu:calU x calU->calU and sigma:calU->calU are defined by mu(U,V):=u\_delta(U bold-dot V) and sigma(U):=u\_delta(U\textasciicircum{}dagger), are global C\textasciicircum{}1 and, for every U,V,Z in calU, mu(J,U)=mu(U,J)=U, sigma(J)=J, ||mu(U,V)-U bold-dot V||<=C\_grp*epsilon\_r, ||sigma(U)-U\textasciicircum{}dagger||<=C\_grp*epsilon\_r, ||mu(mu(U,V),Z)-mu(U,mu(V,Z))||<=C\_grp*epsilon\_r, ||mu(sigma(U),U)-J||<=C\_grp*epsilon\_r, and ||mu(U,sigma(U))-J||<=C\_grp*epsilon\_r; (A\_6) for every U\_0,U\_1 in calU and q in [0,1] satisfying ||U\_1-U\_0||<=q, C\_path*q<=1/4, and C\_path*(q+epsilon\_r*q+q\textasciicircum{}2)<delta\_*-C\_pol*(epsilon\_r*delta\_*+delta\_*\textasciicircum{}2), the path H(-,U\_0,U\_1):[0,1]->calU given by H(t,U\_0,U\_1):=u\_delta((1-t)*U\_0+t*U\_1) is defined, is jointly continuous in its displayed variables, and joins U\_0 to U\_1, and satisfies H(t,cU\_0,cU\_1)=c*H(t,U\_0,U\_1) for every c in U(1) and t in [0,1]; (A\_7) for every s in \{+1,-1\}, set chi\_s:B\_\{r\_iso\}\textasciicircum{}\{icalH\}(0)->calU by chi\_s(A):=sJ bold-dot (J+A+g\_\{sJ\}(A)), let phi\_\{sJ\}\textasciicircum{}par:chi\_s(B\_\{r\_iso\}\textasciicircum{}\{icalH\}(0))->B\_\{r\_iso\}\textasciicircum{}\{icalH\}(0) be its inverse, and set F\_s:B\_\{r\_iso\}\textasciicircum{}\{icalH\}(0)->icalH by F\_s(A):=phi\_\{sJ\}\textasciicircum{}par(sigma(chi\_s(A))); then sigma maps chi\_s(B\_\{r\_iso\}\textasciicircum{}\{icalH\}(0)) into itself and ||D(F\_s-id)(A)+2I||<=C\_der*(epsilon\_r+r\_iso) for every A in B\_\{r\_iso\}\textasciicircum{}\{icalH\}(0); and (R) set q\_*:=C\_grp*epsilon\_r, set r\_-:=delta\_*-C\_pol*(epsilon\_r*delta\_*+delta\_*\textasciicircum{}2), and set eta\_*:=C\_path*(q\_*+epsilon\_r*q\_*+q\_*\textasciicircum{}2); the guards C\_ch*(epsilon\_r+delta\_*)<=kappa\_ch, C\_pol*(epsilon\_r+delta\_*)<=kappa\_pol, q\_*<r\_-, C\_path*q\_*<=1/4, eta\_*<r\_-, C\_der*(epsilon\_r+r\_iso)<=kappa\_der/4<1, and (1+epsilon\_r)*(1+C\_ch*(epsilon\_r+delta\_*))*r\_iso+q\_*<2delta\_* hold; then, for those same u\_delta,h\_delta,mu,sigma,H, there exist r\_bidx>0 depending only on W, a space breve-calU, and maps breve-mu:breve-calU x breve-calU->breve-calU and breve-sigma:breve-calU->breve-calU such that r\_bidx=r\_iso and breve-calU=calU\_e/U(1), set breve-e:=[J], breve-mu([U],[V])=[mu(U,V)], breve-sigma([U])=[sigma(U)], (breve-calU,breve-mu,breve-e) is a connected H-space, breve-sigma is a smooth left inversion, sigma(cU)=conj(c)*sigma(U), ||sigma(U)-U\textasciicircum{}dagger||<=C\_grp*epsilon\_r, breve-calU is a connected compact orientable smooth manifold without boundary of real dimension (dim\_C calX)-1 and is homeomorphic to a finite simplicial complex, breve-e is an isolated fixed point of breve-sigma with local index +1, J and -J are the only sigma-fixed points in their respective ambient r\_bidx-balls, and for every breve-sigma-fixed class breve-V there exist U\_0 in calU\_e and c,a in U(1) such that [U\_0]=breve-V, sigma(U\_0)=c*U\_0, a\textasciicircum{}2=c, sigma(a*U\_0)=a*U\_0, and sigma(-a*U\_0)=-a*U\_0. \steptag{validated} \\[6pt]
\step{1.1} Fix an arbitrary Stage-1 polar witness tuple W, algebra, epsilon\_r, and displayed data satisfying every antecedent of node 1. Universal instantiation of the validated external lem-stage1-bound-inversion-isolation supplies universal witnesses e\_quot\textasciicircum{}r>0 and epsilon\_B\textasciicircum{}r=min\{epsilon\_*\textasciicircum{}r,e\_quot\textasciicircum{}r\}>0 and, for these same u\_delta,h\_delta,mu,sigma,H, witnesses r\_bidx, breve-calU, breve-mu, breve-sigma having every conclusion of that external contract: r\_bidx=r\_iso, the quotient and induced-map formulas, connected H-space and smooth-left-inversion properties, scalar equivariance and norm estimate, the connected compact orientable smooth boundaryless manifold statement with dimension, the isolated quotient fixed point and index +1, and the two ambient isolation statements. \steptag{validated} \\[4pt]
\step{1.2} For the breve-calU supplied in node 1.1, lem-stage1-quotient-finite-cw applies because the ambient algebra is finite-dimensional and exact-unit and node 1.1 states that breve-calU=calU\_e/U(1) is a compact smooth manifold without boundary. Hence breve-calU is homeomorphic to a finite simplicial complex (and consequently has finite CW type). \steptag{validated} \\[4pt]
\step{1.3} For every fixed class breve-V of the induced map breve-sigma supplied in node 1.1, there exist U\_0 in calU\_e and c,a in U(1) such that [U\_0]=breve-V, sigma(U\_0)=c*U\_0, a\textasciicircum{}2=c, sigma(a*U\_0)=a*U\_0, and sigma(-a*U\_0)=-a*U\_0. \steptag{validated} \\[6pt]
\step{1.3.1} Let breve-V be breve-sigma-fixed. Since breve-calU=calU\_e/U(1), choose U\_0 in calU\_e with [U\_0]=breve-V. Using breve-sigma([U])=[sigma(U)] from node 1.1, fixedness gives [sigma(U\_0)]=[U\_0]; by the defining scalar-orbit equivalence of this quotient, there is c in U(1) such that sigma(U\_0)=c*U\_0. \steptag{validated} \\[4pt]
\step{1.3.2} For the phase c in node 1.3.1, there exists a in U(1) with a\textasciicircum{}2=c: write c=exp(i*theta) for some real theta and take a=exp(i*theta/2), which has modulus one and squares to c. \steptag{validated} \\[4pt]
\step{1.3.3} For U\_0,c,a from nodes 1.3.1-1.3.2, scalar equivariance sigma(zU)=conj(z)*sigma(U) from node 1.1 yields sigma(a*U\_0)=conj(a)*c*U\_0 and sigma(-a*U\_0)=-conj(a)*c*U\_0. Since a is in U(1), conj(a)*a=1; with c=a\textasciicircum{}2 this gives conj(a)*c=a. Therefore sigma(a*U\_0)=a*U\_0 and sigma(-a*U\_0)=-a*U\_0, proving all assertions of node 1.3. \steptag{validated} \\[4pt]
\end{proofbox}

\end{document}
